
\subsubsection{Erst umformen, dann die Wurzel}%%$\sqrt{\,\vphantom{b}}$}

Steht eine Gleichung schon in der Form $\left(T(x)\right)^n = B$ da, so ist es
sinnvoll, gleich die Wurzel zu ziehen. Ansonsten versuchen Sie, die
Gleichung in diese Form zu bringen.

$$(x-5)^8 - 9 =0$$
\TNT{8}{
  Potenz-Form:
  $$(x-5)^8 = 9$$
  8-te Wurzel ziehen (Plus / Minus beachten, da gerade Wurzel)

  $$x-5 = \pm\sqrt[8]{9}$$

  $x$ separieren

  $$x = 5 \pm \sqrt[8]{9}$$

  $$\lx = \{5-\sqrt[8]{9}; 5+\sqrt[8]{9}\}$$
}

\begin{rezept}{Potenzgleichung}{}
  Bringen Sie die Potenzgleichungen am besten zuerst in die Form:
  $$\left(T(x)\right)^n = B$$
  Beispiel:
  $$\left(x^2-2x+6\right)^9 = \frac73$$
  \end{rezept}
\newpage



\subsection*{Aufgabenblatt}
%%  OLAT Arbeitsblatt
\GESO{\olatLinkArbeitsblatt{Potenz- und
    Wurzelgleichungen}{https://olat.bms-w.ch/auth/RepositoryEntry/6029794/CourseNode/101937264586296}{Kap. 1.2
Aufg. 2. a) bis 2. f) [optional: Kap. 1.3: Aufg. 3. a) und b)]}}%% END olatLinkArbeitsblatt
\TALS{\olatLinkArbeitsblatt{Potenz- und
    Wurzelgleichungen}{https://olat.bms-w.ch/auth/RepositoryEntry/6029786/CourseNode/105177901239296}{Kap. 1.2
Aufg. 2. a) bis 2. f) [optional: Kap. 1.3: Aufg. 3. a) und b)]}}%% END olatLinkArbeitsblatt

\newpage

\subsection{Spezielle Exponenten}
\subsubsection{negative
  Exponenten}\index{Exponenten!negative}\index{negative Exponenten}
Bei Potenzgleichungen mit negativen Exponenten ist am einfachsten
folgendes Potenzgesetz anzuwenden:

$$a^{-n} = \frac1{a^n}$$

Beachte, dass es manchmal zwei Lösungen, manchmal keine gibt.

Beispiel:

$$x^{-4} = 16$$

\TNT{8}{
  Potenzgesetz:   $$\frac1{x^4} = 16$$
  Kehrwert: $$x^4 = \frac1{16}$$
  4. Wurzel: $$x=\pm \sqrt[4\,\,] {\frac1{16}}= \pm
  \frac1{\sqrt[4]{16}} = \pm \frac12$$
$$\lx = \left\{-\frac12; \frac12\right\}$$}%% END TNT

\subsection*{Aufgaben}

%%  OLAT Arbeitsblatt
\GESO{\olatLinkArbeitsblatt{Potenz- und
    Wurzelgleichungen}{https://olat.bms-w.ch/auth/RepositoryEntry/6029794/CourseNode/101937264586296}{Kap. 2.1
Aufg. 4. a) bis 4. d)}}%% END olatLinkArbeitsblatt
\TALS{\olatLinkArbeitsblatt{Potenz- und
    Wurzelgleichungen}{https://olat.bms-w.ch/auth/RepositoryEntry/6029786/CourseNode/105177901239296}{Kap. 2.1
Aufg. 4. a) bis 4. d)}}%% END olatLinkArbeitsblatt



\newpage

\subsubsection{Rationale Exponenten\GESO{ (optional)}}\index{Exponenten!rationale}\index{Wurzeln}


\TNT{2}{Anstelle von $\sqrt[n]{x}$ hatten wir auch  $x^{\frac{1}{n}}$
  geschrieben.}%% end TNT

Beispiel: $\sqrt[4]{x^3} = 125$

\TNT{10}{

\begin{tabular}{rccl}
%%                & $\sqrt[4]{x^3}$               &=&   125                \\
  $\Rightarrow$ & $x^{\frac{3}{4}}$               &=&   $125$                \\
  $\Rightarrow$ & $(x^{\frac{3}{4}})^\frac{4}{3}$  &=& $(125)^\frac{4}{3}$  \\
  $\Rightarrow$ & $x^{\frac{3}{4}\cdot\frac{4}{3}}$  &=& $(\sqrt[3]{125})^4$  \\
  $\Rightarrow$ & $x^1$                          &=& $5^4$  \\
  $\Rightarrow$ & $x$                            &=& $625$  
\end{tabular}

}%% END TNT

\begin{bemerkung}{Wurzelgleichung}{}
Elementare Wurzelgleichungen\index{Wurzelgleichung} können aufgefasst werden als
Potenzgleichung mit rationalen Exponenten.
\end{bemerkung}

%%%%%%%%%%%%%%%%%%%%%%%%%%%%%%%%%%%%%%%%%%%%%%%%%%%%%%%%%%%%%%%%%%%%%%%


\subsection*{Aufgaben}

%%  OLAT Arbeitsblatt
\GESO{\olatLinkArbeitsblatt{Potenz- und
    Wurzelgleichungen}{https://olat.bms-w.ch/auth/RepositoryEntry/6029794/CourseNode/101937264586296}{Aufg. Kap. 2.2 Aufg. 5}}%% END olatLinkArbeitsblatt
\TALS{\olatLinkArbeitsblatt{Potenz- und
    Wurzelgleichungen}{https://olat.bms-w.ch/auth/RepositoryEntry/6029786/CourseNode/105177901239296}{Aufg. Kap. 2.2 Aufg. 5}}%% END olatLinkArbeitsblatt

\olatLinkGESOKompendium{2.4}{18}{60. bis 65.}
\newpage
